% Generated draft: requires mathematical and source review.
\documentclass[11pt]{article}
\usepackage[margin=25mm]{geometry}
\usepackage{fontspec}
\setmainfont{Latin Modern Roman}
\usepackage{amsmath,amssymb,mathtools,unicode-math}
\setmathfont{Latin Modern Math}
\usepackage{enumitem,needspace,xurl,hyperref}
\setlength{\parindent}{0pt}
\setlength{\parskip}{5pt}
\setlength{\emergencystretch}{3em}
\newcommand{\N}{\mathbb{N}}
\newcommand{\Z}{\mathbb{Z}}
\newcommand{\Q}{\mathbb{Q}}
\newcommand{\R}{\mathbb{R}}
\newcommand{\C}{\mathbb{C}}
\newcommand{\F}{\mathbb{F}}
\newcommand{\A}{\mathbb{A}}
\newcommand{\PP}{\mathbb{P}}
\newcommand{\E}{\mathbb{E}}
\newcommand{\eps}{\varepsilon}
\newcommand{\dd}{\,\mathrm{d}}
\begin{document}
\section*{UnsolvedMath ID: 701429}
\textbf{Source problem code:} catalogue entry
\subsection{Definitions and mathematical statement}
Let \(p\) be any prime, let \(T\) be a finite group of order a power of \(p\), and let \(A\leq \operatorname{Aut}(T)\) satisfy \(A\cong (\Z/p\Z)^2\). Thus \(A\) is an elementary abelian group of automorphisms of order \(p^2\). Form the semidirect product \(P=T\rtimes A\), with multiplication
\[
(t,a)(u,b)=(t\,a(u),ab),
\]
and identify \(T\) with \(T\times\{1_A\}\). Suppose that
\[
\operatorname{ord}_P(x)=p\qquad\text{for every }x\in P\setminus T.
\]
Must it follow that
\[
t^p=1_T\qquad\text{for every }t\in T?
\]
This is the question whether \(T\) has exponent \(p\), where the exponent is the least positive integer \(n\) such that \(t^n=1_T\) for every \(t\in T\). The subgroup condition on \(A\) implies that \(T\) is nontrivial, so the displayed conclusion is equivalent to exponent exactly \(p\). No restriction to odd primes is supplied.
\subsection{Short English statement}
E. Jabara asks whether a finite \(p\)-group must have exponent \(p\) if it admits an elementary abelian automorphism group of order \(p^2\) such that every element outside the original group in the resulting semidirect product has order \(p\).
\subsection{Sources}
\begin{itemize}
\item Local UnsolvedMath registry, supplied source record: permanent ID \(701429\), exact source problem code \texttt{catalogue entry}, title \emph{Kourovka Notebook Problem 17.128}. Upstream revision: \texttt{d371ee917ac9647fbb06cd6fc529deec6f3dda82}.
\item \emph{The Kourovka Notebook (Unsolved Problems in Group Theory)}, 17th issue (2010), Problem 17.128; proposed by E. Jabara. The supplied background also identifies the 21st issue edition (2026), \texttt{arXiv:1401.0300v48}. The supplied collection description names E. I. Khukhro and V. D. Mazurov as editors.
\item The registry labels the problem \emph{open}. This is a supplied status, not an independently verified literature assessment; the cited editions have not been independently consulted for this draft.
\end{itemize}
\end{document}
